\section{An isolated pair: collision time and the surrounding zeros}
\label{sec:analytic-isolated-collision}

\begin{theorem}[NN-A023: an isolated-pair collision with its first interaction correction]
Let $f$ be a real entire function of order at most one, with
only real zeros, and let $b-\delta/2,b+\delta/2$ be adjacent
simple zeros, $\delta>0$. Put
\[
 G(z)=\frac{f(z)}{(z-b)^2-\delta^2/4},\qquad
 L=\frac{G'}G(b)\in\mathbb R,\qquad
 S=\sum_{G(\xi)=0}\frac1{(b-\xi)^2},\qquad
                       \kappa=\delta^2S.
\]
Zeros in $S$ are counted with multiplicity. Assume the
order-one Hadamard product for $G$; its possible linear
exponential factor is allowed and is retained in $L$.
The case $S=0$ is included.

For all sufficiently small $\kappa$, with an absolute
threshold, the two zeros continued from this pair under
$f_{-a}=e^{a\partial_z^2}f$ undergo a unique double collision
on the local time scale $0\le a\le\delta^2/4$. Its time and
real position satisfy
\begin{align}
 a_c&=\delta^2\left(\frac18+\frac{\kappa}{32}
                                      +O(\kappa^2)\right),
                        \label{eq:analytic-isolated-time}\\
 z_c&=b-2a_cL+O(\delta\kappa^{3/2}).
                        \label{eq:analytic-isolated-center}
\end{align}
The $O$ constants are absolute. Before the collision the
tracked zeros are real, and after it they are a nonreal
conjugate pair. The local neighbourhood moves with the real
center $b-2aL$; no smallness assumption on $\delta L$ is needed.

Under A this applies to every adjacent simple pair of $H_0$.
In particular any sequence of such arithmetic pairs with
$\kappa\to0$ has actual collision times
$t_c=-\delta^2(1+\kappa/4+O(\kappa^2))/8$.
This is a conditional statement about the stated pair
parameter, not a proof of how often these pairs occur.
\end{theorem}

\begin{proof}
Hadamard factorization has no exponential term of degree two
or higher, by the order assumption. Therefore
\[
 R(h)=\frac{G(b+h)}{G(b)e^{Lh}}
       =\prod_{G(\xi)=0}
         \left(1+\frac h{b-\xi}\right)
                           \exp\left(-\frac h{b-\xi}\right)
\]
converges locally uniformly, and the proof of
Lemma~\ref{lem:analytic-product-remainder} gives
$|R(h)|\le e^{S|h|^2/2}$. Its first Taylor coefficients are
\begin{equation}\label{eq:analytic-pair-R-expansion}
 R(h)=1-\frac S2h^2+\frac{T_3}{3}h^3+\mathcal E_4(h),
          \qquad T_3=\sum_{G(\xi)=0}(b-\xi)^{-3},
          \qquad |T_3|\le S^{3/2}.
\end{equation}
There is an absolute constant $C$ such that
\begin{equation}\label{eq:analytic-pair-R-fourth}
 |\mathcal E_4(h)|\le
               CS^2|h|^4e^{2S|h|^2}\qquad(h\in\mathbb C).
\end{equation}
Indeed, after putting $v=h\sqrt S$, Cauchy's estimate on
$|v|=1$ bounds the remainder by $2e^{1/2}|v|^4$ for
$|v|\le1/2$. For $|v|\ge1/2$, use the global growth bound,
the explicit quadratic coefficient and
$|T_3|/(3S^{3/2})\le1/3$; each of the four terms is bounded
by a fixed multiple of $|v|^4e^{2|v|^2}$.
If $S=0$, the product is empty and $R=1$, so the claimed
collision and formulas are exact.

Write $a=\delta^2\tau$ and use the moving coordinate
\[
                         z=b-2aL+\delta w.
\]
Completing the Gaussian square, or conjugating the heat
operator by $e^{L(z-b)}$, gives the exact identity
\begin{equation}\label{eq:analytic-pair-scaled-family}
 \frac{f_{-\delta^2\tau}(b-2\delta^2\tau L+\delta w)}
      {\delta^2G(b)\exp(\delta wL-\delta^2\tau L^2)}
  =J(\tau,w)
  =\mathbb E\left[(Y^2-\tfrac14)R(\delta Y)\right],
                \qquad Y=w+\sqrt{2\tau}\,V,\quad V\sim N(0,1).
\end{equation}
For real $\tau\ge0$ this is an ordinary Gaussian convolution
with a real shift $2aL$. Both sides extend holomorphically
in $\tau$ to a complex neighbourhood of $[0,1/4]$:
one may choose a square-root branch locally away from zero;
the expectation is even in that square root and extends
across zero. The order-one growth and Gaussian decay
justify the continuation. The denominator is nonzero.

For bounded complex $w,\tau$, the growth bound for $R$
is dominated by $\exp(C\kappa(|V|^2+1))$.
Consequently \eqref{eq:analytic-pair-R-fourth}, integrated
against the Gaussian, gives a uniform holomorphic expansion
with an $O(\kappa^2)$ remainder:
\begin{align}
 J(\tau,w)
  &=w^2+2\tau-\tfrac14-\frac{\kappa}{2}P_4(\tau,w)
                  +\eta_3 P_5(\tau,w)+O(\kappa^2),
                    \label{eq:analytic-pair-J-expansion}\\
 \eta_3&=\delta^3T_3/3,\qquad |\eta_3|\le\kappa^{3/2}/3,\nonumber\\
 P_4(\tau,w)
  &=w^4+(12\tau-\tfrac14)w^2+12\tau^2-\tfrac12\tau,\nonumber\\
 P_5(\tau,w)
  &=w^5+(20\tau-\tfrac14)w^3
                         +(60\tau^2-\tfrac32\tau)w.\nonumber
\end{align}
All derivatives on smaller compact sets have the same
remainder order by Cauchy's formula. In particular the
quadratic first correction is even in $w$, while the
cubic correction is odd. Uniform constants depend only
on the fixed compact sets, not on $f,\delta,L,S$.

For $0\le\tau\le1/4$, the limit polynomial
$w^2+2\tau-1/4$ has exactly two roots in $|w|<2$ and
none on the boundary. Rouch\'e's theorem therefore tracks
exactly two roots of $J$ in this moving disc for small
$\kappa$, starting with the given pair at $\tau=0$.
Its derivative is $2w+O(\kappa)$ uniformly there and has
a single zero near zero. The implicit function theorem
gives this zero as a real analytic function $w_*(\tau)$
for real $\tau$, with
$w_*(\tau)=O(\kappa^{3/2})$: evaluate
$\partial_wJ(\tau,0)$ using the parity in
\eqref{eq:analytic-pair-J-expansion}, and use
$\partial_w^2J=2+O(\kappa)$.

Set $D(\tau)=J(\tau,w_*(\tau))$. The same expansion yields
\[
 D(\tau)=2\tau-\tfrac14
              -\frac{\kappa}{2}(12\tau^2-\tfrac12\tau)
              +O(\kappa^2),\qquad D'(\tau)=2+O(\kappa).
\]
It has one real zero $\tau_c$ near $1/8$. First
$\tau_c=1/8+O(\kappa)$; evaluating the correction at $1/8$,
where $12\tau^2-\tau/2=1/8$, then gives
\[
                       \tau_c=\frac18+\frac{\kappa}{32}
                                                  +O(\kappa^2).
\]
The vanishing conditions $J=J_w=0$, together with
$J_{ww}=2+O(\kappa)\ne0$, give a double zero, not a higher
collision. Since $D'$ and $J_{ww}$ are positive for real
arguments here, its two branches are real for
$\tau<\tau_c$ and nonreal conjugate for $\tau>\tau_c$.
No other local collision is possible on $[0,1/4]$:
the derivative has only the one critical point and $D$
has only the one zero. Returning to $a,z$ proves
\eqref{eq:analytic-isolated-time} and
\eqref{eq:analytic-isolated-center}. This proof also
accounts for the full pair throughout its local interval;
it does not merely locate a nonreal zero at one time.
\end{proof}

\paragraph{Relation to the classical Lehmer criterion.}
Csordas--Smith--Varga, \emph{Lehmer Pairs of Zeros, the
de Bruijn--Newman Constant $\Lambda$, and the Riemann
Hypothesis}, Constructive Approximation~10 (1994), 107--129,
Definition~1, equations~(1.10)--(1.12), printed pp.~109--110,
uses the endpoint sum
\[
 g=\sum_{\xi\ne b-\delta/2,b+\delta/2}
       \left\{(b-\delta/2-\xi)^{-2}
                       +(b+\delta/2-\xi)^{-2}\right\}.
\]
Under A this is positive. For $\kappa\to0$,
$g=2S+O(\delta^2S^2)$, by expanding the two summands
symmetrically and using $\sum(b-\xi)^{-4}\le S^2$.
Thus $\delta^2g=2\kappa+O(\kappa^2)$, and the pairs above
eventually satisfy their $\delta^2g<4/5$ criterion.
Their Theorem~1, equations~(1.13)--(1.14), printed p.~110,
states the lower bound
\[
 \frac{(1-\tfrac54\delta^2g)^{4/5}-1}{8g}\le\Lambda.
\]
It does not assert that this displayed lower bound is an
exact collision time. Its first two expansion terms agree
with \eqref{eq:analytic-isolated-time} after changing the
sign of time. The argument here identifies an actual
unique local double collision and its interaction correction.
The original printed pages were inspected directly.

\paragraph{Geometric and representation meaning.}
The chosen two-point divisor determines a finite local
rank-two zero family in the moving coordinate. Its
discriminant has a simple zero at $\tau_c$, so after a
quadratic base change the two branches are separated.
For the resulting finite locally free rank-two algebra,
the exact \'{e}tale/discriminant criterion is SGA~1,
Expos\'e~I, section~4, Proposition~10; section~3,
Proposition~1 gives the equivalent vanishing of relative
differentials. This criterion is applied to the actual
local finite algebra, not to all of the infinite divisor.
The first correction is the inverse-square trace of the
complementary zero divisor, while its logarithmic first
trace produces only the displayed real translation.
This is a concrete local discriminant computation; it
does not make the complementary analytic divisor into a
finite pure Weil object. In the Toeplitz/character
dictionary of NN-A020--A022, each such collision creates
one conjugate pair of complex product parameters.
The relevant exterior-power operations have the locator
SGA~6, Expos\'e~V, Example~2.2(b) and Lemma~2.2.1.
No integral canonical-basis statement follows from
the analytic inverse-square parameter.
